\section{Introduction}\label{sec:introduction}

The ionization energy of an atom measures the cost of removing electrons
while leaving its nucleus fixed. The first few electrons removed belong to
the outer part of the atom, whereas the leading total energy is largely
determined by the much denser inner region. Thomas--Fermi theory, introduced
by Thomas and Fermi \cite{Thomas1927,Fermi1927}, describes the leading
large-charge energy. Lieb and Simon established this connection to the
many-body Schr\"odinger theory \cite{LiebSimon1973,LiebSimon1977}.
Their energy theorem \cite[Theorem~III.1]{LiebSimon1977} concerns terms of
order $Z^{7/3}$. Subtracting two such total-energy asymptotics does not
control a much smaller ionization energy: the remainder may exceed the
difference being sought.

The generalized ionization conjecture asserts that Thomas--Fermi theory
nevertheless determines the leading cost of removing many outer electrons,
provided their number remains small compared with the nuclear charge.
Solovej formulated this statement through upper and lower large-$Z$ limits,
followed by a large removed-electron-number limit
\cite[Section~3, Equation~(1)]{Solovej2016}. We prove its energy part for
the full two-spin Coulomb Hamiltonian, in the stronger joint regime in which
both the number removed and the ratio of nuclear charge to that number
tend to infinity.

Earlier full-model estimates cover different ionization regimes. Seco,
Sigal and Solovej bounded the first ionization energy of a neutral atom
by $CZ^{20/21}$, using the improvement recorded in their note added in
proof \cite[Theorem~1 and p.~315]{SecoSigalSolovej1990}.
Ivrii's estimates imply agreement with the magnitude of the Thomas--Fermi
chemical potential for successive one-electron ionization energies of positive
atoms when $Z^{5/7}=o(Z-N)$ and $Z-N=o(Z)$, where $N$ is the number of
electrons \cite[Theorem~6.2.1]{Ivrii2017}. These results leave open the
joint neutral-to-positive-ion energy asymptotic addressed here.

Uniform excess-charge bounds were proved in
Thomas--Fermi--von Weizs\"acker theory by Benguria and Lieb
\cite{BenguriaLieb1985}; Solovej proved uniform ionization bounds in reduced
Hartree--Fock theory \cite{Solovej1991} and then in unrestricted
Hartree--Fock theory \cite{Solovej2003}. The Thomas--Fermi--von Weizs\"acker
and Hartree--Fock results concern different variational problems. The passage to
arbitrary correlated fermionic states requires additional screening
estimates, obtained here through conditional many-body states.

\subsection{The theorem and its normalization}

In atomic units, let $Z$ be a positive integer and let
\begin{equation}\label{eq:hamiltonian}
 H_{Z,N}
 =\sum_{i=1}^N\left(-\frac12\Delta_i-\frac{Z}{|x_i|}\right)
   +\sum_{1\le i<l\le N}\frac1{|x_i-x_l|}
\end{equation}
act on $\bigwedge^N L^2(\R^3;\C^2)$. Write $E_Z(N)=\inf\sigma(H_{Z,N})$
and $E_Z(0)=0$. For integers $Z>m\ge1$, define
\[
 I_m(Z)=E_Z(Z-m)-E_Z(Z).
\]
The kinetic coefficient of the corresponding Thomas--Fermi functional is
$T=\frac3{10}(3\pi^2)^{2/3}$. More explicitly, put
\[
 \mathcal T_Z(\rho)
 =T\int_{\R^3}\rho^{5/3}
 -Z\int_{\R^3}\frac{\rho(x)}{|x|}\,dx
 +\frac12\iint_{\R^3\times\R^3}
          \frac{\rho(x)\rho(y)}{|x-y|}\,dx\,dy,
\]
and let $E_Z^{\TF}(N)$ be its infimum over nonnegative densities of mass $N$
and finite displayed terms. Set
$I_m^{\TF}(Z)=E_Z^{\TF}(Z-m)-E_Z^{\TF}(Z)$.
The constant in the conjecture is characterized by the classical
Thomas--Fermi weak-ionization law
\cite[Section~6, Corollary~5]{BenilanBrezis2004}:
\begin{equation}\label{eq:tf-constant}
 \lim_{Z\to\infty}I_m^{\TF}(Z)=a_{\TF}m^{7/3},\qquad m>0.
\end{equation}
Our identification of the coefficient in \cref{sec:fixed-sectors}
recovers this characterization in the present normalization.

\begin{theorem}[Generalized ionization energy]\label{thm:main}
For full nonrelativistic Coulomb atoms with two electron spin states,
\begin{equation}\label{eq:joint-limit}
 \frac{I_m(Z)}{m^{7/3}}\longrightarrow a_{\TF}>0
 \quad\text{whenever}\quad m\longrightarrow\infty,
 \qquad \frac Zm\longrightarrow\infty,
\end{equation}
where $Z>m\ge1$ are integers. In particular,
\begin{align}
 \lim_{m\to\infty}
  \frac{\limsup_{Z\to\infty} I_m(Z)}{m^{7/3}}
   &=a_{\TF},\label{eq:iterated-upper}\\
 \lim_{m\to\infty}
  \frac{\liminf_{Z\to\infty} I_m(Z)}{m^{7/3}}
   &=a_{\TF}.\label{eq:iterated-lower}
\end{align}
The inner limits are taken over integer $Z$ with $m$ fixed. No convergence
of $I_m(Z)$ at a fixed $m$ is asserted or assumed.
\end{theorem}

The limits in \cref{eq:iterated-upper,eq:iterated-lower} are precisely the
energy part of the generalized ionization conjecture formulated by
Solovej \cite[Section~3, Equation~(1)]{Solovej2016}. The joint assertion \cref{eq:joint-limit} strengthens that order of limits.
For historical comparison, Solovej established uniform first-ionization
upper bounds and sharp generalized radius estimates in unrestricted Hartree--Fock
theory \cite[Theorems~3.8 and~1.5]{Solovej2003}; the generalized conjecture
and its Hartree--Fock history are discussed in \cite[Section~3]{Solovej2016}.
Here the energy is the exact spectral infimum of
\cref{eq:hamiltonian}, with no restriction to Slater determinants.

For this same full two-spin Coulomb model, the companion
\cite[Corollary~1.3]{Foundation} also proves
$0<c\le E_Z(Z-1)-E_Z(Z)\le C<\infty$ uniformly for every integer $Z\ge1$.
This first-ionization bound does not assert convergence at fixed $m$ and
is distinct from the large-$m$ limits above.

\subsection{A price per electron}

The proof first allows the particle number to vary. For fixed $Z$, abbreviate
$E_n=E_Z(n)$, and introduce the unconstrained energy and the priced minimum
\[
 E=\inf_{n\ge0}E_n,\qquad
 P(\lambda)=\min_{n\ge0}\bigl(E_n+\lambda n\bigr),\qquad\lambda>0.
\]
The minimum exists, as verified in \cref{sec:pricing}. Its advantage is
local: replacing the exterior electrons by a trial state need not preserve
their number. The price accounts for the change exactly. At the length
scale $s=\lambda^{-1/4}$, the removed charge is of order $s^{-3}$.
The main intermediate result identifies its coefficient for every minimizing
sector, including at prices where several sectors minimize.

\begin{theorem}[Deficit at a prescribed price]\label{thm:priced-deficit}
There is a universal number $q>0$ such that, for every sequence satisfying
\[
 \lambda=s^{-4},\qquad s\longrightarrow0,
 \qquad Zs^3\longrightarrow\infty,
\]
and every choice of $n$ minimizing $E_n+\lambda n$, one has
\[
 s^3(Z-n)\longrightarrow q.
\]
The number $q$ is the exterior Coulomb charge of the unique classical solution, tending uniformly to zero at infinity,
of
\[
 \Delta F=4\pi k\bigl((F-1)_+\bigr)^{3/2}
 \quad\text{on }\R^3\setminus\{0\},
 \qquad k=(5T/3)^{-3/2},
\]
that is bounded above and below by positive multiples of $|x|^{-4}$ near
zero. The density $k((F-1)_+)^{3/2}$ has bounded support, and
$F(x)=q/|x|$ outside that support.
Moreover,
\[
 a_{\TF}=\frac37q^{-4/3}.
\]
\end{theorem}

\subsection{Analytic precedents and proof structure}

The local kinetic comparison combines two classical constructions.
Coherent wave packets provide the upper semiclassical energy
\cite[Section~V.A.1]{Lieb1981ThomasFermi}; see also
\cite[Lemma~8.2]{Solovej2003}. Neumann cube eigenvalues and their lattice
count supply a lower comparison with the same leading coefficient
\cite[Theorems~III.10--III.13]{LiebSimon1977}. Their remainders are controlled
by local particle counts. This is different from the coarse
Lieb--Thirring kinetic inequality \cite{LiebThirring1975}, whose positive
universal constant is sufficient for preliminary density bounds; the
spectral splitting proof used in the companion follows Rumin
\cite[Theorem~1.4 and Section~3.1]{Rumin2011}.

Screened exterior potentials and outward Thomas--Fermi comparisons are
central to Solovej's Hartree--Fock argument
\cite[Sections~6, 10, and~12]{Solovej2003}. Our use of conditional states
requires a separate transfer from local quantum energies to conditional
fields, while accounting for changes in particle number. The eventual
singular profile belongs to the classical Sommerfeld comparison theory
\cite[Section~IV]{LiebSimon1977}, including the nonspherical comparisons
and chemical-potential estimates of
\cite[Lemma~4.4, Remark~4.5, Theorem~4.6, and Corollary~4.7]{Solovej2003}.
The issue here is to obtain that profile from the correlated quantum
state and preserve its positive singularity through conditioning.


We use the localization, screening, and Thomas--Fermi estimates proved in
the companion \emph{Uniform excess charge for Coulomb molecules and the
outer radius of neutral atoms} \cite{Foundation}.
In particular, its uniform neutral-sector bound $E_Z(Z)-E\le C$ fixes
the zero-price integration constant, and its local-field estimates retain
explicit control of an arbitrary energy offset
\cite[\FNeutralOffset\ and \FLocalField]{Foundation}.
The required positive-price conditional comparisons and barrier
propagation are proved in \cref{sec:comparison,sec:barriers}.
Their offset dependence is essential: a minimizing priced state has
energy at most $E+Cs^{-7}$, so the allowed offset grows as $s\to0$.
A fixed-offset conclusion from the neutral-atom theory cannot by itself
yield \cref{thm:main}.

\Cref{sec:pricing} obtains the rough deficit bound and proves that the
rescaled exterior electron mass vanishes beyond a fixed radius. In
\cref{sec:comparison}, noisy observations of the electrons permit local
conditional comparisons while controlling the cost of conditioning. A
Thomas--Fermi trial state in a small ball is compared with the priced
many-body state. A fresh spatial cut records the local particles and leaves
a core outside the patch, whose conditional screened potential is harmonic
on the patch. Matching this core comparison to the noisy-observation law
yields simultaneous upper and lower implications between local density
and local potential.
Two posterior laws occur in this argument; keeping them distinct is
necessary for the Coulomb comparison.

The remaining difficulty is to retain the singular inner charge through
successive averaging. \Cref{sec:barriers} constructs a positive subsolution
at the innermost scale and propagates it through the observation scales.
The construction pays for exceptional conditional data with an error
density of vanishing expected mass. The weak differential inequality is
proved on open neighborhoods of the retained branches before taking their
maximum, so no derivative is restricted to an activity interface.
\Cref{sec:limit} then selects conditional data, rescales by $s$, and obtains
a decaying solution of the unit-price Thomas--Fermi equation. The barrier
prevents loss of the singularity. A maximum-principle argument identifies
its Sommerfeld coefficient and proves uniqueness, which forces the same
deficit for every minimizing sector.

Finally, \cref{sec:fixed-sectors} integrates the slopes of the priced
minimum and brackets the desired particle number between two minimizing
sectors. Only monotonicity of $E_Z(N)$ is used in this step; convexity in
$N$ is not needed. The analogous Thomas--Fermi calculation identifies the
coefficient with \cref{eq:tf-constant}. A sequential argument then gives
the specified order of limits in
\cref{eq:iterated-upper,eq:iterated-lower}.
